\documentclass{article}
\usepackage{fontspec}
\setmainfont{Libertinus Serif}
\usepackage[a4paper, margin=1in]{geometry}
\usepackage{amsmath, amsthm, amssymb, mathtools, enumitem, hyperref}
\usepackage[portuguese]{babel}

\setlist[description]{leftmargin=1cm,labelindent=1cm}

\theoremstyle{plain}
\newtheorem{thm}{Teorema}[section]
\newtheorem{prop}[thm]{Proposição}
\newtheorem{lem}[thm]{Lema}
\newtheorem{cor}[thm]{Corolário}

\theoremstyle{definition}
\newtheorem{defn}{Definição}[section]
\newtheorem{exmp}[defn]{Exemplo}
\newtheorem{prblm}[defn]{Problema}
\newtheorem{prcp}[defn]{Princípio}

\theoremstyle{remark}
\newtheorem*{rem}{Nota}

\DeclareMathOperator{\im}{im}
\DeclareMathOperator{\spn}{span}

\geometry{
 a4paper,
 left=3cm,
 top=2cm,
 right=2cm,
 bottom=2cm
}

\begin{document}
\begin{center}
    \textbf{\large UNIVERSIDADE FEDERAL DO RIO GRANDE DO SUL} \\
    \textbf{\large INSTITUTO DE MATEMÁTICA E ESTATÍSTICA} \\
    \vspace{3cm}
    \textbf{\Large ÁLGEBRA LINEAR II} \\
    \vspace{1cm}
    \textbf{\large LISTA DE DEFINIÇÕES E RESULTADOS} \\
    \vspace{3cm}

    \textbf{LEONARDO SUAREZ}

    \vfill
    \textbf{\large Porto Alegre, RS} \\
    \textbf{\large \today}
\end{center}

\newpage
\tableofcontents
\newpage

\section{Espaços Vetoriais}

\subsection{\(\mathbb{R}^n\) e \(\mathbb{C}^n\)}
\begin{prop} (propriedades da aritmética complexa).
    \begin{enumerate}
        \item[\textup{(a)}] \textup{\textbf{comutatividade}} \(\alpha + \beta = \beta + \alpha\) e \(\alpha \beta = \beta \alpha \quad \forall \alpha, \beta \in \mathbb{C}\)
        \item[\textup{(b)}] \textup{\textbf{associatividade}} \((\alpha + \beta) + \gamma = \alpha + (\beta + \gamma)\) e \((\alpha \beta) \gamma = \alpha (\beta \gamma) \quad \forall \alpha, \beta, \gamma \in \mathbb{C}\)
        \item[\textup{(c)}] \textup{\textbf{identidades}} \(\alpha + 0 = \alpha\) e \(\alpha 1 = \alpha \quad \forall \alpha \in \mathbb{C}\)
        \item[\textup{(d)}] \textup{\textbf{inverso aditivo}} \(\forall \alpha \in \mathbb{C}\ \exists!\beta \in \mathbb{C}: \alpha + \beta = 0\).
        \item[\textup{(e)}] \textup{\textbf{inverso multiplicativo}} \(\forall \alpha \in \mathbb{C}, \alpha \neq 0\ \exists! \beta \in \mathbb{C}: \alpha \beta = 1\).
        \item[\textup{(f)}] \textup{\textbf{distributividade}} \(\alpha(\beta + \gamma) = \alpha \beta + \alpha \gamma \quad \forall \alpha, \beta, \gamma \in \mathbb{C}\)
    \end{enumerate}
\end{prop}

\begin{defn}
    Suponha que \(n\) seja um inteiro não-negativo. Uma \emph{lista} de \emph{tamanho} \(n\) é uma coleção ordenada de \(n\) elementos (que podem ser números, outras listas, ou outros objetos mais abstratos). Duas listas são iguais sse têm o mesmo tamanho e os mesmos elementos na mesma ordem.
\end{defn}

\begin{defn}
    A \emph{soma} em \(\mathbb{F}^n\) é definida somando as coordenadas correspondentes:
    \[
        (x_1, \ldots, x_n) + (y_1, \ldots, y_n) = (x_1 + y_1, \ldots, x_n + y_n).
    \]
\end{defn}

\begin{prop} (comutatividade da soma em \(\mathbb{F}^n\)).
    Para todo \(x, y \in \mathbb{F}^n\), \(x + y = y + x\).
\end{prop}

\begin{defn}
    Para todo \(x \in \mathbb{F}^n\), o \emph{inverso aditivo} de \(x\) é o vetor \(-x \in \mathbb{F}^n\) tal que \(x + (-x) = 0\). Explicitamente, se \(x = (x_1, \ldots, x_n)\), então \(-x = (-x_1, \ldots, -x_n)\).
\end{defn}

\begin{defn}
    O \emph{produto} de um número \(\lambda\) e um vetor em \(\mathbb{F}^n\) é computado multiplicando cada coordenada do vetor por \(\lambda\):
    \[
        \lambda (x_1, \ldots, x_n) = (\lambda x_1, \ldots, \lambda x_n);
    \]
    \noindent onde \(\lambda \in \mathbb{F}\) e \((x_1, \ldots, x_n) \in \mathbb{F}^n\).
\end{defn}

\begin{defn}
    Uma \emph{soma} em um conjunto \(V\) é uma função que associa um elemento \(u + v \in V\) a cada par de elementos \(u, v \in V\). Um \emph{produto por escalar} em \(V\) é uma função que associa um elemento \(\lambda v \in V\) a cada escalar \(\lambda \in \mathbb{F}\) e cada elemento \(v \in V\). 
\end{defn}

\subsection{Definição de Espaço Vetorial}
\begin{defn}
    Um \emph{espaço vetorial} é um conjunto \(V\) munido com as operações de soma e produto por escalar que satisfaz as seguintes propriedades.
    \begin{enumerate}
        \item[(A1)] \textbf{comutatividade} \(u + v = v + u\), para todo \(u, v \in V\).
        \item[(A2)] \textbf{associatividade} \((u + v) + w = u + (v + w)\), para todo \(u, v, w \in V\).
        \item[(A3)] \textbf{identidade aditiva} Existe \(0 \in V\) tal que \(v + 0 = v\), para todo \(v \in V\).
        \item[(A4)] \textbf{inverso aditivo} Para todo \(v \in V\) existe \(-v \in V\) tal que \(v + (-v) = 0\).
        \item[(A5)] \textbf{identidade multiplicativa} \(1v = v\), para todo \(v \in V\), onde \(1\) é o elemento identidade em \(\mathbb{F}\).
        \item[(A6)] \textbf{distributividade} \(a(u + v) = au + av\) e \((a+b)v=av + bv\), para todo \(a, b \in \mathbb{F}\) e todo \(u, v \in V\).
    \end{enumerate}
\end{defn}

\begin{prop}
    Um espaço vetorial possui identidade aditiva única.
\end{prop}

\begin{prop}
    Todo elemento de um espaço vetorial possui inverso aditivo único.
\end{prop}

\begin{prop}
    \(0v = 0\) e \(a0 = 0\), para todo \(a \in \mathbb{F}\) e todo \(v \in V\).
\end{prop}

\begin{prop}
    \((-1)v = -v\), para todo \(v \in V\).
\end{prop}

\subsection{Subespaços}
\begin{defn}
    Um subconjunto \(U\) de um espaço vetorial \(V\) é dito um \emph{subespaço} de \(V\) se \(U\) é um espaço vetorial com as mesmas operações de soma e produto por escalar definidas em \(V\).
\end{defn}

\begin{prop} (condições para subespaço).
    Um subconjunto \(U\) de um espaço vetorial \(V\) é um subespaço de \(V\) sse os seguintes três itens são satisfeitos.
    \begin{itemize}
        \item \(0 \in U\).
        \item Se \(u, v \in U\), então \(u + v \in U\).
        \item Se \(u \in U\) e \(a \in \mathbb{F}\), então \(au \in U\).
    \end{itemize}
\end{prop}

\begin{defn}
    Suponha que \(V_1,\ldots,V_m\) são subespaços de \(V\). A \emph{soma} \(V_1 + \cdots + V_m\) é o conjunto de todas as somas possíveis de elementos de \(V_1, \ldots, V_m\), isto é,
    \[
        V_1 + \cdots + V_m = \{v_1 + \cdots + v_m : v_1 \in V_1, \ldots, v_m \in V_m\}.
    \]
\end{defn}

\begin{prop}
    Se \(V_1, \ldots, V_m\) são subespaços de \(V\), então \(V_1 + \cdots + V_m\) é o menor subespaço de \(V\) que contém \(V_1, \ldots, V_m\).
\end{prop}

\begin{defn}
    Sejam \(V_1, \ldots, V_m\) subespaços de \(V\). A soma \(V_1 + \cdots + V_m\) é dita \emph{direta} se todo \(v \in V_1 + \cdots + V_m\) pode ser escrito de maneira única como \(v = v_1 + \cdots + v_m\), onde \(v_k \in V_k\). Se a soma é direta, escrevemos \(V = V_1 \oplus \cdots \oplus V_m\).
\end{defn}

\begin{prop} (condição para soma direta).
    Suponha que \(V_1, \ldots, V_m\) são subespaços de \(V\). Então \(V = V_1 \oplus \cdots \oplus V_m\) sse a única maneira de escrever \(0\) como \(0 = v_1 + \cdots + v_m\), onde \(v_k \in V_k\) para cada \(k\), é com \(v_1 = \cdots = v_m = 0\).
\end{prop}

\begin{prop}
    Sejam \(U\) e \(W\) subespaços de \(V\). Então \[U + W \ \text{é soma direta} \iff U \cap W = \{0\}.\]
\end{prop}

\section{Espaços Vetoriais de Dimensão Finita}

\subsection{Span e Independência Linear}
\begin{prop} (span é o menor subespaço que contém todos os vetores).
    O span de uma lista de vetores em \(V\) é o menor subespaço de \(V\) que contém todos os vetores da lista.
\end{prop}

\begin{defn}
    Um espaço vetorial \(V\) é dito de \emph{dimensão finita} se existe uma lista finita de vetores em \(V\) que gera \(V\). Caso contrário, \(V\) é dito de \emph{dimensão infinita}.
\end{defn}

\begin{defn}
    Uma lista \(v_1, \ldots, v_n\) de vetores em \(V\) é dita \emph{linearmente independente} se a única escolha de escalares \(a_1, \ldots, a_n\) em \(\mathbb{F}\) que satisfaz
    \begin{equation}
        a_1 v_1 + a_2 v_2 + \cdots + a_n v_n = 0
    \end{equation}
    \noindent é \(a_1 = a_2 = \cdots = a_n = 0\). Caso contrário, a lista é dita \emph{linearmente dependente}. A lista vazia é dita linearmente independente.
\end{defn}

\begin{lem} (lema de dependência linear).
    Suponha que \(v_1, \ldots, v_n\) é linearmente dependente em \(V\). Então existe \(j \in \{1, \ldots, n\}\) tal que
    \begin{equation}
        v_j \in \spn(v_1, \ldots, v_{j-1}).
    \end{equation}
    \noindent Ainda, se \(j\) satisfaz a condição acima e o \(j\)-ésimo vetor for removido da lista, então o span dos vetores restantes é igual a \(\spn(v_1, \ldots, v_n)\).
\end{lem}

\begin{prop}
    Em um espaço vetorial \(V\) de dimensão finita, o tamanho de qualquer lista linearmente independente é menor ou igual ao tamanho de qualquer lista geradora.
\end{prop}

\begin{prop}
    Todo subespaço de um espaço vetorial de dimensão finita é de dimensão finita.
\end{prop}

\subsection{Bases}
\begin{defn}
    Uma \emph{base} de um espaço vetorial \(V\) é uma lista de vetores \(v_1, \ldots, v_n\) em \(V\) que é linearmente independente e gera \(V\).
\end{defn}

\begin{prop} (critério para base).
    Uma lista de vetores \(v_1, \ldots, v_n\) em um \(V\) é uma base de \(V\) sse todo vetor em \(V\) pode ser escrito de maneira única como uma combinação linear de \(v_1, \ldots, v_n\), isto é,
    \begin{equation}
        v = a_1 v_1 + a_2 v_2 + \cdots + a_n v_n,
    \end{equation}
    \noindent onde \(a_1, \ldots, a_n\) são escalares em \(\mathbb{F}\).
\end{prop}

\begin{prop}
    Toda lista geradora em um espaço vetorial pode ser reduzida a uma base do espaço vetorial.
\end{prop}

\begin{cor}
    Todo espaço vetorial de dimensão finita possui uma base.
\end{cor}

\begin{prop}
    Toda lista linearmente independente em um espaço vetorial de dimensão finita pode ser estendida a uma base do espaço vetorial.
\end{prop}

\begin{prop}
    Suponha que \(V\) seja um espaço vetorial de dimensão finita e \(U\) é um subespaço de \(V\). Então existe um subespaço \(W\) de \(V\) tal que \(V = U \oplus W\).
\end{prop}

\subsection{Dimensão}
\begin{prop} (o tamanho da base não depende da base).
    Se \(V\) é um espaço vetorial de dimensão finita, então todas as bases de \(V\) têm o mesmo tamanho.
\end{prop}

\begin{defn}
    A \emph{dimensão} de um espaço vetorial \(V\) de dimensão finita é o tamanho de qualquer base de \(V\).
\end{defn}

\begin{prop}
    Se \(V\) é um espaço vetorial de dimensão finita e \(U\) é um subespaço de \(V\), então \(\dim(U) \leq \dim(V)\).
\end{prop}

\begin{prop}
    Se \(V\) é um espaço vetorial de dimensão finita, então toda lista linearmente independente de vetores em \(V\) com tamanho igual a \(\dim(V)\) é uma base de \(V\).
\end{prop}

\begin{cor}
    Se \(V\) é um espaço vetorial de dimensão finita, então todo subespaço \(U\) de \(V\) com \(\dim(U) = \dim(V)\) é igual a \(V\).
\end{cor}

\begin{prop}
    Se \(V\) é um espaço vetorial de dimensão finita, então toda lista geradora de vetores em \(V\) com tamanho igual a \(\dim(V)\) é uma base de \(V\).
\end{prop}

\begin{prop} (dimensão de soma).
    Se \(V_1\) e \(V_2\) são subespaços de um espaço vetorial de dimensão finita \(V\), então
    \begin{equation}
        \dim(V_1 + V_2) = \dim(V_1) + \dim(V_2) - \dim(V_1 \cap V_2).
    \end{equation}
\end{prop}

\section{Transformações Lineares}

\subsection{Espaços Vetoriais de Transformaçõess Lineares}
\begin{defn}
    Uma transformação linear de \(V\) em \(W\) é uma função \(T: V \to W\) com as seguintes propriedades.
    \begin{enumerate}
        \item \(T(u + v) = Tu + Tv\), para todo \(u, v \in V\).\hfill\emph{(aditividade)}
        \item \(T(\lambda u) = \lambda Tu\), para todo \(u \in V\) e todo escalar \(\lambda \in \mathbb{F}\).\hfill\emph{(homogeneidade)}
    \end{enumerate}
\end{defn}

\begin{lem}
    Suponha que \(v_1,\ldots,v_n\) é uma base de \(V\) e \(w_1,\ldots,w_n \in W\). Então existe uma única transformação linear \(T: V \to W\) tal que \(Tv_k = w_k\), para \(k = 1, \ldots, n\).
\end{lem}

\begin{defn}
    Suponha \(S, T \in \mathcal{L}(V, W)\) e \(\lambda \in \mathbb{F}\). A \emph{soma} \(S + T\) e o \emph{produto por escalar} \(\lambda T\) são as transformações lineares de \(V\) em \(W\) definidas por
    \[
        (S + T)(v) = S(v) + Tv \quad \text{e} \quad (\lambda T)(v) = \lambda (Tv),
    \]
    \noindent para todo \(v \in V\).
\end{defn}

\begin{cor}
    Com as operações de soma e produto escalar definidas acima, \(\mathcal{L}(V, W)\) é um espaço vetorial sobre \(\mathbb{F}\).
\end{cor}

\begin{defn}
    Se \(T \in \mathcal{L}(U, V)\) e \(S \in \mathcal{L}(V, W)\), então o \emph{produto} \(ST\) é a transformação linear de \(U\) em \(W\) definida por
    \[
        (ST)(u) = S(Tu),
    \]
    \noindent para todo \(u \in U\).
\end{defn}

\begin{prop} (propriedades algébricas do produto de transformações lineares).
    \begin{itemize}
        \item \(\left(T_1 T_2\right)T_3 = T_1\left(T_2T_3\right)\) sempre que \(T_1, T_2, T_3\) são transformações lineares tais que os produtos são definidos.\hfill\emph{(associatividade)}
        \item \(T_1I = IT_1 = T_1\) sempre que \(T_1 \in \mathcal{L}(V, W)\), onde o primeiro \(I\) é a identidade em \(V\) e o segundo \(I\) é a identidade em \(W\).\hfill\emph{(elemento neutro)}
        \item \(\left(S_1 + S_2\right)T = S_1T + S_2T\) e \(S\left(T_1 + T_2\right) = ST_1 + ST_2\), sempre que \(T, T_1, T_2 \in \mathcal{L}(U, V)\) e \(S, S_1, S_2 \in \mathcal{L}(V, W)\).\hfill\emph{(distributividade)}
    \end{itemize}
\end{prop}

\begin{cor}
    Suponha que \(T\) seja uma transformação linear de \(V\) em \(W\). Então \(T(0) = 0\).
\end{cor}

\subsection{Espaço Nulo e Imagem}
\begin{defn}
    Para \(T \in \mathcal{L}(V, W)\), o \emph{espaço nulo} de \(T\) é o subconjunto de \(V\) definido por
    \[
        \ker(T) = \{v \in V : Tv = 0\}.
    \]
\end{defn}

\begin{prop}
    Seja \(T \in \mathcal{L}(V, W)\). Então \(\ker(T)\) é um subespaço de \(V\).
\end{prop}

\begin{defn}
    Uma função \(T: V \to W\) é \emph{injetora} se \(Tv_1 = TTv_2\) implica \(v_1 = v_2\), para todo \(v_1, v_2 \in V\).
\end{defn}

\begin{prop}
    Seja \(T \in \mathcal{L}(V, W)\). Então \(T\) é injetora se e somente se \(\ker(T) = \{0\}\).
\end{prop}

\begin{defn}
    Para \(T \in \mathcal{L}(V, W)\), a \emph{imagem} de \(T\) é o subconjunto de \(W\) definido por
    \[
        \text{Im}(T) = \{Tv : v \in V\}.
    \]
\end{defn}

\begin{prop}
    Seja \(T \in \mathcal{L}(V, W)\). Então \(\im(T)\) é um subespaço de \(W\).
\end{prop}

\begin{defn}
    Uma função \(T: V \to W\) é \emph{sobrejetora} se para todo \(w \in W\) existe \(v \in V\) tal que \(Tv = w\).
\end{defn}

\begin{thm} (teorema fundamental das transformações lineares).
    Suponha que \(V\) seja um espaço vetorial de dimensão finita e que \(T \in \mathcal{L}(V, W)\). Então \(\im(T)\) é de dimensão finita e
    \begin{equation}
        \dim V = \dim \ker T + \dim \im T.
    \end{equation}
\end{thm}

\begin{proof}
    Sejam \(u_1, \ldots, u_n\) uma base de \(\ker T\) e \(u_1, \ldots, u_n, v_1, \ldots, v_m\) uma base de \(V\). Então temos que \(\dim \ker T = n\) e \(\dim V = n+m\). Resta provar que \(\dim \im T = m\).
    Escreve
    \[
        v = a_1 u_1 + \cdots + a_n u_n + b_1 v_1 + \cdots + b_m v_m
    \]
    \noindent para qualquer \(v \in V\) e com \(\{a_i\}_{i=1}^n, \{b_j\}_{j=1}^m \in \mathbb{F}\). Aplica \(T\) e fica 
    \[
        Tv = b_1 Tv_1 + \cdots + b_m Tv_m ,
    \]
    \noindent isto é, \(Tv_1, \ldots, Tv_m\) gera \(\im T\). Para provar que \(Tv_1, \ldots, Tv_m\) é LI, considere a equação
    \begin{gather}
        c_1 Tv_1 + \cdots + c_m Tv_m = 0 \\
        T(c_1 v_1 + \cdots + c_m v_m) = 0,
    \end{gather}
    \noindent com \(\{c_j\}_{j=1}^m \in \mathbb{F}\), e nota que \(c_1 v_1 + \cdots + c_m v_m \in \ker T\). Então escreve 
    \begin{gather}
        c_1 v_1 + \cdots + c_m v_m = d_1 u_1 + \cdots + d_n u_n \\
        c_1 v_1 + \cdots + c_m v_m - d_1 u_1 - \cdots - d_n u_n = 0.
    \end{gather}
    Como \(u_1, \ldots, u_n, v_1, \ldots, v_m\) é base, a última equação acima possui apenas a solução trivial \(c_1 = \cdots = c_m = d_1 = \cdots = d_n = 0\). Logo a lista \(Tv_1, \ldots, Tv_m\) é LI e forma, portanto, uma base para \(\im T\).
\end{proof}

\begin{cor}
    Suponha que \(V\) e \(W\) são espaços vetoriais de dimensão finita com \(\dim V > \dim W\). Então nenhuma transformação linear de \(V\) em \(W\) é injetora.
\end{cor}

\begin{cor}
    Suponha que \(V\) e \(W\) são espaços vetoriais de dimensão finita com \(\dim V < \dim W\). Então nenhuma transformação linear de \(V\) em \(W\) é sobrejetora.
\end{cor}

\subsection{Matrizes}
\begin{defn}
    Suponha que \(m\) e \(n\) são inteiros não-negativos. Uma \emph{matriz} \(m\) por \(n\) é uma tabela retangular de elementos de \(\mathbb{F}\) com \(m\) linhas e \(n\) colunas:
    \[
        A = \begin{bmatrix}
            A_{11} & A_{12} & \cdots & A_{1n} \\
            A_{21} & A_{22} & \cdots & A_{2n} \\
            \vdots & \vdots & \ddots & \vdots \\
            A_{m1} & A_{m2} & \cdots & A_{mn}
        \end{bmatrix},
    \]
    \noindent onde \(A_{jk} \in \mathbb{F}\) é o elemento na \(j\)-ésima linha e \(k\)-ésima coluna de \(A\).
\end{defn}

\begin{defn} (matriz de uma transformação linear).
    Suponha que \(T \in \mathcal{L}(V, W)\) e que \(v_1, \ldots, v_n\) é uma base de \(V\) e \(w_1, \ldots, w_m\) é uma base de \(W\). A \emph{matriz} de \(T\) com respeito a essas bases é a matriz \(\mathcal{M}(T)\) de tamanho \(m\) por \(n\) cujas entradas \(A_{jk}\) são definidas por
    \[
        Tv_k = A_{1k} w_1 + \cdots + A_{mk} w_m.
    \]    
\end{defn}

\begin{defn}
    A \emph{soma de duas matrizes de mesmo tamanho} é a matriz obtida somando as entradas correspondentes:
    \[
        A + B = \begin{bmatrix}
            A_{11} + B_{11} & A_{12} + B_{12} & \cdots & A_{1n} + B_{1n} \\
            A_{21} + B_{21} & A_{22} + B_{22} & \cdots & A_{2n} + B_{2n} \\
            \vdots & \vdots & \ddots & \vdots \\
            A_{m1} + B_{m1} & A_{m2} + B_{m2} & \cdots & A_{mn} + B_{mn}
        \end{bmatrix}.
    \]
\end{defn}

\begin{prop} (matriz da soma de transformações lineares).
    Sejam \(S, T \in \mathcal{L}(V, W)\). Então \(\mathcal{M}(S + T) = \mathcal{M}(S) + \mathcal{M}(T)\).
\end{prop}

\begin{defn}
    O produto de uma matriz por um escalar é a matriz obtida multiplicando cada entrada da matriz pelo escalar:
    \[
        \lambda A = \begin{bmatrix}
            \lambda A_{11} & \lambda A_{12} & \cdots & \lambda A_{1n} \\
            \lambda A_{21} & \lambda A_{22} & \cdots & \lambda A_{2n} \\
            \vdots & \vdots & \ddots & \vdots \\
            \lambda A_{m1} & \lambda A_{m2} & \cdots & \lambda A_{mn}
        \end{bmatrix}.
    \]
\end{defn}

\begin{prop} (matriz do produto por escalar de uma transformação linear).
    Sejam \(T \in \mathcal{L}(V, W)\) e \(\lambda \in \mathbb{F}\). Então \(\mathcal{M}(\lambda T) = \lambda \mathcal{M}(T)\).
\end{prop}

\begin{prop} \((\textup{dim } \mathbb{F}^{m,n} = mn)\).
    Suponha que \(m\) e \(n\) sejam inteiros positivos. Com as operações de soma e produto por escalar definidas acima, \(\mathbb{F}^{m,n}\) é um espaço vetorial sobre \(\mathbb{F}\) com dimensão \(mn\).    
\end{prop}

\begin{defn} (multiplicação de matrizes).
    Sejam \(A\) uma matriz \(m\) por \(n\) e \(B\) uma matriz \(n\) por \(p\). O \emph{produto} \(AB\) é a matriz \(m\) por \(p\) cujas entradas são definidas por
    \[
        (AB)_{jk} = \sum_{r=1}^{n} A_{jr} B_{rk}.
    \]
\end{defn}

\begin{prop} (matriz do produto de transformações lineares).
    Sejam \(T \in \mathcal{L}(U, V)\) e \(S \in \mathcal{L}(V, W)\). Então \(\mathcal{M}(ST) = \mathcal{M}(S) \mathcal{M}(T)\).
\end{prop}

\begin{prop} (entrada do produto de matrizes é igual ao produto de vetores linha por coluna).
    Sejam \(A\) uma matriz \(m\) por \(n\) e \(B\) uma matriz \(n\) por \(p\). Então a entrada \((j, k)\) de \(AB\) é igual ao produto do \(j\)-ésimo vetor linha de \(A\) pelo \(k\)-ésimo vetor coluna de \(B\):
    \[
        (AB)_{jk} = A_{j.} B_{.k},
    \]
    \noindent se \(1 \leq j \leq m\) e \(1 \leq k \leq p\).
\end{prop}

\begin{prop} (coluna do produto de matrizes é igual ao produto da matriz pela coluna).
    Sejam \(A\) uma matriz \(m\) por \(n\) e \(B\) uma matriz \(n\) por \(p\). Então a \(k\)-ésima coluna de \(AB\) é igual ao produto de \(A\) pela \(k\)-ésima coluna de \(B\):
    \[
        (AB)_{.k} = A B_{.k},
    \]
    \noindent se \(1 \leq k \leq p\).
\end{prop}

\begin{cor} (combinação linear de colunas).
    Sejam \(A\) uma matriz \(m\) por \(n\) e \(b\) uma matriz \(n\) por \(1\). Então o produto \(Ab\) é uma combinação linear das colunas de \(A\) com os coeficientes dados pelas entradas de \(b\):
    \[
        Ab = b_1 A_{.1} + b_2 A_{.2} + \cdots + b_n A_{.n}.
    \]
\end{cor}

\begin{cor} (multiplicação de matrizes como combinação linear de colunas).
    Suponha que \(C\) é uma matriz \(m\) por \(c\) e \(R\) é uma matriz \(c\) por \(n\).
    \begin{enumerate}
        \item[\textup{(a)}] Se \(k \in \{1, \ldots, n\}\), então a \(k\)-ésima coluna de \(CR\) é uma combinação linear das colunas de \(C\) com os coeficientes dados pelas entradas da \(k\)-ésima coluna de \(R\).
        \item[\textup{(b)}] Se \(j \in \{1, \ldots, m\}\), então a \(j\)-ésima linha de \(CR\) é uma combinação linear das linhas de \(R\) com os coeficientes dados pelas entradas da \(j\)-ésima linha de \(C\).
    \end{enumerate}
\end{cor}

\begin{defn}
    Seja \(A\) uma matriz \(m\) por \(n\). O \emph{posto} de \(A\), denotado por \(\text{rank}(A)\), é a dimensão do span das colunas de \(A\) em \(\mathbb{F}^{m1}\). O \emph{posto transposto} de \(A\), denotado por \(\text{rank}(A^t)\), é a dimensão do span das linhas de \(A\) em \(\mathbb{F}^{1n}\).
\end{defn}

\begin{defn}
    A transposta de uma matriz \(A\), denotada por \(A^t\), é a matriz obtida trocando as linhas por colunas, isto é, a entrada \((j, k)\) de \(A^t\) é a entrada \((k, j)\) de \(A\):
    \[
        (A^t)_{jk} = A_{kj}.
    \]
\end{defn}

\begin{prop} (fatoração coluna-linha).
    Seja \(A\) uma matriz \(m\) por \(n\) com entradas em \(\mathbb{F}\) e posto \(c \geq 1\). Então existem uma matriz \(C\) \(m\) por \(c\) e uma matriz \(R\) \(c\) por \(n\), ambas com entradas em \(\mathbb{F}\), tais que \(A = CR\).
\end{prop}

\begin{prop} (posto é igual ao posto transposto).
    Se \(A\) é uma matriz com entradas em \(\mathbb{F}\), então \(\textup{rank}(A) = \textup{rank}(A^t)\).
\end{prop}

\subsection{Invertibilidade e Isomorfismos}
\begin{defn}
    Uma transformação linear \(T \in \mathcal{L}(V, W)\) é dita \emph{invertível} se existe uma transformação linear \(S \in \mathcal{L}(W, V)\) tal que \(ST\) é igual ao operador identidade em \(V\) e \(TS\) é igual ao operador identidade em \(W\). Uma transformação linear \(S \in \mathcal{L}(W, V)\) que satisfaz \(ST = I\) e \(TS = I\) é dita a \emph{inversa} de \(T\).
\end{defn}

\begin{prop}
    Uma transformação linear invertível tem uma única inversa. Se \(T \in \mathcal{L}(V, W)\) é invertível, então \(T^{-1} \in \mathcal{L}(W, V)\) é o único elemento de \(\mathcal{L}(W, V)\) tal que \(T^{-1}T = I\) e \(TT^{-1} = I\).
\end{prop}

\begin{prop}
    Uma transformação linear é invertível sse é bijetora.
\end{prop}

\begin{prop}
    Sejam \(V\) e \(W\) espaços vetoriais de dimensão finita e que \(\dim V = \dim W\). Então uma transformação linear \(T \in \mathcal{L}(V, W)\) é invertível sse \(T\) é injetora sse \(T\) é sobrejetora.
\end{prop}

\begin{proof}
    Sejam \(V\) e \(W\) espaços vetoriais de dimensão finita e \(\dim V = \dim W\). Seja \(T \in \mathcal{L}(V, W)\). Suponha que \(T\) é injetora. Do Teorema do Posto, temos 
    \begin{gather}
        \dim V = \dim \ker T + \dim \im T \\
        \dim W = \dim \im T
    \end{gather}
    Logo, \(T\) é sobrejetora. Agora, suponha que \(T\) é sobrejetora. Novamente, do Teorema do Posto, escreve 
    \begin{gather}
        \dim V = \dim \ker T + \dim W \\
        \dim \ker T = 0,
    \end{gather} 
    \noindent o que prova que \(T\) é injetora. Portanto, temos a equivalência 
    \[
        T \text{ é invertível } \iff T \text{ é injetora } \iff T \text{ é sobrejetora}.
    \]
\end{proof}

\begin{prop}
    Sejam \(V\) e \(W\) espaços vetoriais de dimensão finita, \(\text{dim } V = \text{dim } W\), \(S \in \mathcal{L}(V, W)\) e \(T \in \mathcal{L}(W, U)\). Então \(ST = I\) sse \(TS = I\).
\end{prop}

\begin{defn}
    Um \emph{isomorfismo} é uma transformação linear invertível. Se existe um isomorfismo de \(V\) em \(W\), então \(V\) e \(W\) são ditos \emph{isomorfos}.
\end{defn}

\begin{prop}
    Dois espaços vetoriais de dimensão finita são isomorfos sse têm a mesma dimensão.
\end{prop}

\begin{proof}
    Sejam \(V\) e \(W\) espaços vetoriais de dimensão finita isomorfos. Então existe uma transformação linear \(T\) invertível de \(V\) em \(W\). Do Teorema do Núcleo e da Imagem, temos 
    \[
        \dim V = \dim \ker T + \dim \im T.
    \]
    Como \(T\) é injetora e sobrejetora, vem
    \[
        \dim V = \dim W,
    \]
    \noindent o que prova a primeira direção. Suponha agora que \(\dim V = \dim W\). Sejam \(v_1, \ldots, v_n\) uma base de \(V\) e \(w_1, \ldots, w_n\) uma base de \(W\). Define \(T \in \mathcal{L}(V, W)\) como 
    \[
        T(a_1 v_1 + \cdots + a_n v_n) = a_1 w_1 + \cdots + a_n w_n.
    \]
    Dado que \(v_1, \ldots, v_n\) é base de \(V\), \(T\) está bem-definida. Nota que \(a_1 w_1 + \cdots + a_n w_n = 0\) possui apenas a solução trivial, logo \(T\) é injetora. Por fim, como \(\dim V = \dim W\), \(T\) é também sobrejetora.
\end{proof}

\begin{prop}
    Suponha \(v_1, \ldots, v_n\) é uma base de \(V\) e \(w_1, \ldots, w_n\) é base de \(W\). Então \(\mathcal{M}\) é um isomorfismo de \(\mathcal{L}(V, W)\) em \(\mathbb{F}^{m,n}\).
\end{prop}

\begin{cor}
    Sejam \(V\) e \(W\) espaços vetoriais de dimensão finita. Então \(\mathcal{L}(V, W)\) é de dimensão finita e
    \[
        \dim \mathcal{L}(V, W) = (\dim V)(\dim W).
    \]
\end{cor}

\begin{defn} (matriz de um vetor).
    Seja \(v \in V\) e \(v_1,\ldots,v_n\) uma base de \(V\). A \emph{matriz de \(v\)} com respeito a essa base é a matriz \(n\) por \(1\) definida por
    \[
        \mathcal{M}(v) = \begin{bmatrix}
            b_1 \\
            b_2 \\
            \vdots \\
            b_n
        \end{bmatrix},
    \]
    \noindent onde \(v = b_1 v_1 + b_2 v_2 + \cdots + b_n v_n\).
\end{defn}

\begin{cor}
    Suponha \(T \in \mathcal{L}(V, W)\) e \(v_1,\ldots,v_n\) é uma base de \(V\) e \(w_1,\ldots,w_m\) é uma base de \(W\). Seja \(1 \leq k \leq n\). Então a \(k\)-ésima coluna de \(\mathcal{M}(T)\) é igual a \(\mathcal{M}(Tv_k)\).
\end{cor}

\begin{prop} (transformações lineares atuam como multiplicação de matrizes).
    Suponha \(T \in \mathcal{L}(V, W)\), \(v_1,\ldots,v_n\) é uma base de \(V\) e \(w_1,\ldots,w_m\) é uma base de \(W\). Se \(v \in V\), então
    \[
        \mathcal{M}(Tv) = \mathcal{M}(T) \mathcal{M}(v).
    \]
\end{prop}

\begin{prop}
    Sejam \(V\) e \(W\) espaços vetoriais de dimensão finita e \(T \in \mathcal{L}(V, W)\). Então \(\dim \im(T)\) é igual ao posto de \(\mathcal{M}(T)\).
\end{prop}

\subsection{Dualidade}

\begin{defn}
    Um \emph{funcional linear} em \(V\) é uma transformação linear de \(V\) em \(\mathbb{F}\). Em outras palavras, um funcional linear é um elemento de \(\mathcal{L}(V, \mathbb{F})\).
\end{defn}

\begin{defn}
    O \emph{dual} de \(V\), denotado por \(V'\), é o espaço vetorial de todos os funcionais lineares sobre \(V\), ou seja, \(V' = \mathcal{L}(V, \mathbb{F})\).
\end{defn}

\begin{prop}
    Se \(V\) é um espaço vetorial de dimensão finita, então \(V'\) é de dimensão finita e \(\dim V' = \dim V\).
\end{prop}

\begin{defn}
    Se \(v_1, \ldots, v_n\) é uma base de \(V\), então a \emph{base dual} de \(V'\) com respeito a essa base é a lista \(\varphi_1, \ldots, \varphi_n\) de funcionais lineares em \(V\) definida por
    \[
        \varphi_i(v_j) = \begin{cases}
            1 & \text{se } i = j \\
            0 & \text{se } i \neq j
        \end{cases}, \quad \forall i,j = 1,\ldots,n.
    \]
\end{defn}

\section{Autovalores e Autovetores}
\subsection{Subespaços Invariantes}

\begin{defn}
    Uma transformação linear de um espaço vetorial em si mesmo é dita um \emph{operador} no espaço vetorial. O conjunto de todos os operadores em \(V\) é denotado por \(\mathcal{L}(V)\).    
\end{defn}

\begin{defn}
    Suponha \(T \in \mathcal{L}(V)\). Um subespaço \(U\) de \(V\) é dito \emph{invariante} sob \(T\) se \(Tu \in U\) para todo \(u \in U\).
\end{defn}

\begin{defn}
    Suponha \(T \in \mathcal{L}(V)\). Um escalar \(\lambda \in \mathbb{F}\) é dito um \emph{autovalor} de \(T\) se existe \(v \in V\) tal que \(v \neq 0\) e \(Tv = \lambda v\). 
\end{defn}

\begin{prop} (condições equivalentes para autovalor).
    Seja \(V\) um espaço vetorial de dimensão finita, \(T \in \mathcal{L}(V)\) e \(\lambda \in \mathbb{F}\). Então as seguintes condições são equivalentes.
    \begin{enumerate}
        \item[\textup{(a)}] \(\lambda\) é um autovalor de \(T\).
        \item[\textup{(b)}] \(T - \lambda I\) não é injetora.
        \item[\textup{(c)}] \(T - \lambda I\) não é sobrejetora.
        \item[\textup{(d)}] \(T - \lambda I\) não é invertível.
    \end{enumerate}
\end{prop}

\begin{defn}
    Suponha \(T \in \mathcal{L}(V)\) e \(\lambda\) é um autovalor de \(T\). O \emph{autovetor} de \(T\) correspondente a \(\lambda\) é qualquer vetor \(v \in V\) tal que \(v \neq 0\) e \(Tv = \lambda v\).
\end{defn}

\begin{thm}
    Seja \(T \in \mathcal{L}(V)\). Então toda lista de autovetores de \(T\) correspondente a autovalores distintos de \(T\) é linearmente independente.
\end{thm}

\begin{proof}
    Seja \(\{\lambda_i\}_{i=1}^p\) uma lista de autovalores distintos de \(T\). Suponha, por absurdo, que a lista de autovetores \(\{v_i\}_{i=1}^p\) seja LD. Tome a menor lista LD \(\{v_i\}_{i=1}^n\), \(1 \leq n \leq p\) tal que todos os coeficientes da equação 
    \[
        a_1 v_1 + \cdots + a_n v_n = 0
    \]
    \noindent sejam não nulos. Aplica \(T - \lambda_n I\) e fica 
    \[
        a_1 (\lambda_1 - \lambda_n) + \cdots + a_{n-1} (\lambda_{n-1} - \lambda_n)v_{n-1} = 0.
    \]
    Como os autovalores são distintos, todos os coeficientes são não nulos. Então, encontramos uma lista de autovetores LD de tamanho \(n-1\), o que contradiz a minimalidade de \(\{v_i\}_{i=1}^n\).
\end{proof}

\begin{defn}
    Suponha \(T \in \mathcal{L}(V)\) e \(m\) é um inteiro positivo.
    \begin{itemize}
        \item \(T^m \in \mathcal{L}(V)\) é definido por \(T^m = T T \cdots T\) (\(m\) vezes).
        \item \(T^0\) é definido como o operador identidade em \(V\).
        \item Se \(T\) é invertível com inversa \(T^{-1}\), então \(T^{-m} = (T^{-1})^m\).
    \end{itemize}
\end{defn}

\begin{defn}
    Suponha \(T \in \mathcal{L}(V)\) e \(p \in \mathcal{P}(\mathbb{F})\) é um polinômio dado por
    \[    p(z) = a_0 + a_1 z + a_2 z^2 + \cdots + a_m z^m.\]
    \noindent para todo \(z \in \mathbb{F}\). Então \(p(T) \in \mathcal{L}(V)\) é o operador em \(V\) definido por
    \[    p(T) = a_0 I + a_1 T + a_2 T^2 + \cdots + a_m T^m.\]
\end{defn}

\begin{defn} (produto de polinômios).
    Se \(p\), \(q \in \mathcal{P}(\mathbb{F})\), então o \emph{produto} \(pq\) é o polinômio em \(\mathcal{P}(\mathbb{F})\) definido por
    \[    (pq)(z) = p(z) q(z),\]
    \noindent para todo \(z \in \mathbb{F}\).
\end{defn}

\begin{prop} (propriedades multiplicativas).
    Sejam \(p, q \in \mathcal{P}(\mathbb{F})\) e \(T \in \mathcal{L}(V)\). Então
    \begin{enumerate}
        \item[(a)] \((pq)(T) = p(T) q(T)\);
        \item[(b)] \(p(T) q(T) = q(T) p(T)\).
    \end{enumerate}
\end{prop}

\begin{prop} (espaço nulo e imagem de \(p(T)\) são invariantes sob \(T\)).
    Sejam \(T \in \mathcal{L}(V)\) e \(p \in \mathcal{P}(\mathbb{F})\). Então os subespaços \(\ker p(T)\) e \(\im p(T)\) são invariantes sob \(T\).
\end{prop}

\subsection{Matrizes Triangular-Superiores}

\begin{defn}
    Seja \(T \in \mathcal{L}(V)\). A \emph{matriz} de \(T\) com respeito a uma base \(v_1, \ldots, v_n\) de \(V\) é a matriz \(\mathcal{M}(T)\) de tamanho \(n\) por \(n\)
    \[
    \mathcal{M}(T) = \begin{pmatrix}
        A_{11} & A_{12} & \cdots & A_{1n} \\
        A_{21} & A_{22} & \cdots & A_{2n} \\
        \vdots & \vdots & \ddots & \vdots \\
        A_{n1} & A_{n2} & \cdots & A_{nn}
    \end{pmatrix},
    \]
    \noindent cujas entradas \(A_{jk}\) são definidas por
    \[
        Tv_k = A_{1k} v_1 + A_{2k} v_2 + \cdots + A_{nk} v_n.
    \]
    A notação \(\mathcal{M}(T, (v_1, \ldots, v_n))\) também é usada para denotar a matriz de \(T\) com respeito à base \(v_1, \ldots, v_n\).
\end{defn}

\begin{defn}
    A \emph{diagonal} de uma matriz quadrada consiste nas entradas na linha do canto superior esquerdo até a linha do canto inferior direito.
\end{defn}

\begin{defn}
    Uma matriz quadrada é dita \emph{triangular-superior} se todas as entradas abaixo da diagonal são zero.
\end{defn}

\begin{prop} (condições para matriz triangular-superior).
    Seja \(T \in \mathcal{L}(V)\) e \(v_1, \ldots, v_n\) uma base de \(V\). Então as seguintes condições são equivalentes.
    \begin{enumerate}
        \item[\textup{(a)}] A matriz de \(T\) com respeito à base \(v_1, \ldots, v_n\) é triangular-superior.
        \item[\textup{(b)}] \(\spn(v_1, \ldots, v_k)\) é invariante sob \(T\), para \(k = 1, \ldots, n\).
        \item[\textup{(c)}] \(Tv_k \in \spn(v_1, \ldots, v_k)\), para \(k = 1, \ldots, n\).
    \end{enumerate}
\end{prop}

\begin{prop}
    Suponha que \(T \in \mathcal{L}(V)\) e que \(V\) tem uma base tal que a matriz de \(T\) com respeito a essa base é triangular-superior com entradas diagonais \(\lambda_1, \ldots, \lambda_n\). então
    \[
        (T - \lambda_1I)\cdots(T - \lambda_n I) = 0.
    \]
\end{prop}

\begin{prop} (determinação de autovalores a partir de matriz triangular-superior).
    Suponha que \(T \in \mathcal{L}(V)\) tem uma matriz triangular-superior com respeito a alguma base de \(V\). Então os autovalores de \(T\) são exatamente as entradas na diagonal da matriz.
\end{prop}

\begin{prop} (condição necessária e suficiente para existência de matriz triangular-superior).
    Suponha que \(V\) é um espaço vetorial de dimensão finita e \(T \in \mathcal{L}(V)\). Então \(T\) tem uma matriz triangular-superior com respeito a alguma base de \(V\) sse o polinômio mínimo de \(T\) é igual a \((z - \lambda_1) \cdots (z - \lambda_m)\), para alguns \(\lambda_1, \ldots, \lambda_m \in \mathbb{F}\).
\end{prop}

\begin{cor} (se \(\mathbb{F}=\mathbb{C}\), então todo operador em \(V\) tem matriz triangular-superior).
    Sejam \(V\) um espaço vetorial complexo de dimensão finita e \(T \in \mathcal{L}(V)\). Então \(T\) tem uma matriz triangular-superior com respeito a alguma base de \(V\).
\end{cor}

\subsection{Operadores Diagonalizáveis}

\begin{defn}
    Uma \emph{matriz diagonal} é uma matriz quadrada cujas entradas fora da diagonal são todas zero.
\end{defn}

\begin{defn}
    Um operador em \(V\) é dito \emph{diagonalizável} se o operador tem uma matriz diagonal com respeito a alguma base de \(V\).
\end{defn}

\begin{defn}
    Suponha que \(T \in \mathcal{L}(V)\) e \(\lambda \in \mathbb{F}\). O \emph{autoespaço} de \(T\) correspondente a \(\lambda\) é o subespaço \(E(\lambda, T)\) de \(V\) definido por
    \[
        E(\lambda, T) = \ker(T - \lambda I) = \{v \in V : Tv = \lambda v\}.
    \]
    \noindent Então, \(E(\lambda, T)\) é o conjunto de todos os autovetores de \(T\) correspondentes a \(\lambda\), juntamente com o vetor nulo.
\end{defn}

\begin{prop} (soma de autoespaços é direta).
    Suponha que \(T \in \mathcal{L}(V)\) e que \(\lambda_1, \ldots, \lambda_m\) são autovalores distintos de \(T\). Então a soma \(E(\lambda_1, T) + \cdots + E(\lambda_m, T)\) é direta. Além disso, se \(V\) é de dimensão finita, então
    \[
        \dim E(\lambda_1, T) + \cdots + \dim E(\lambda_m, T) \leq \dim V.
    \]
\end{prop}

\begin{thm} (condições para diagonalizabilidade).
    Suponha que \(V\) é um espaço vetorial de dimensão finita e \(T \in \mathcal{L}(V)\). Sejam \(\lambda_1,\ldots,\lambda_m\) autovalores distintos de \(T\). Então as seguintes condições são equivalentes.
    \begin{enumerate}
        \item[\textup{(a)}] \(T\) é diagonalizável.
        \item[\textup{(b)}] \(V\) tem uma base de autovetores de \(T\).
        \item[\textup{(c)}] \(V = E(\lambda_1, T) \oplus \cdots \oplus E(\lambda_m, T)\).
        \item[\textup{(d)}] \(\dim V = \dim E(\lambda_1, T) + \cdots + \dim E(\lambda_m, T)\).
    \end{enumerate}
\end{thm}

\begin{prop} (autovalores suficientes implicam diagonalizabilidade).
    Suponha que \(V\) é um espaço vetorial de dimensão finita e \(T \in \mathcal{L}(V)\). Se \(T\) tem \(\dim V\) autovalores distintos, então \(T\) é diagonalizável.
\end{prop}

\section{Espaços com Produto Interno}
\subsection{Produtos Internos e Normas}

\begin{defn}
    Para \(x, y \in \mathbb{R}^n\), o \emph{produto interno} de \(x\) e \(y\) é o número real \(x \cdot y\) definido por
    \[
        x \cdot y = x_1 y_1 + \cdots + x_n y_n,
    \]
    \noindent onde \(x = (x_1, \ldots, x_n)\) e \(y = (y_1, \ldots, y_n)\).
\end{defn}

\begin{defn}
    Um \emph{produto interno} em \(V\) é uma função que toma um par ordenado \((u, v)\) de elementos de \(V\) e atribui a ele um número escalar \(\langle u, v \rangle \in \mathbb{F}\) e tem as seguintes propriedades.
    \begin{itemize}
        \item \textbf{positividade:} \(\langle v, v \rangle \geq 0\), para todo \(v \in V\), e \(\langle v, v \rangle = 0\) sse \(v = 0\);
        \item \textbf{aditividade na primeira variável:} \(\langle u + v, w \rangle = \langle u, w \rangle + \langle v, w \rangle\), para todos \(u, v, w \in V\);
        \item \textbf{homogeneidade na primeira variável:} \(\langle \lambda v, w \rangle = \lambda \langle v, w \rangle\), para todo \(\lambda \in \mathbb{F}\) e todos \(v, w \in V\);
        \item \textbf{conjugação simétrica:} \(\langle w, v \rangle = \overline{\langle v, w \rangle}\), para todos \(v, w \in V\).
    \end{itemize}
\end{defn}

\begin{defn}
    Um \emph{espaço com produto interno} é um par \((V, \langle \cdot, \cdot \rangle)\), onde \(V\) é um espaço vetorial sobre \(\mathbb{F}\) e \(\langle \cdot, \cdot \rangle\) é um produto interno em \(V\).
\end{defn}

\begin{prop} (propriedades básicas do produto interno).
    \begin{enumerate}
        \item [\textup{(a)}] Para cada \(v \in V\) fixado, a função que toma \(u \in V\) em \(\langle u, v \rangle\) é uma transformação linear em \(V\).
        \item [\textup{(b)}] \(\langle 0, v \rangle = 0\) para todo \(v \in V\).
        \item [\textup{(c)}] \(\langle v, 0 \rangle = 0\) para todo \(v \in V\).
        \item [\textup{(d)}] \(\langle u, v + w \rangle = \langle u, v \rangle + \langle u, w \rangle\), para todos \(u, v, w \in V\).
        \item [\textup{(e)}] \(\langle u, \lambda v \rangle = \overline{\lambda} \langle u, v \rangle\), para todo \(\lambda \in \mathbb{F}\) e todos \(u, v \in V\).
    \end{enumerate}
\end{prop}

\begin{defn}
    Para \(v \in V\), a \emph{norma} de \(v\), denotada por \(\|v\|\), é o número real não-negativo definido por \(\|v\| = \sqrt{\langle v, v \rangle}\).
\end{defn}

\begin{prop} (propriedades básicas da norma).
    Seja \(v \in V\).
    \begin{enumerate}
        \item [\textup{(a)}] \(\|v\| \geq 0\) e \(\|v\| = 0\) sse \(v = 0\).
        \item [\textup{(b)}] \(\|\lambda v\| = |\lambda| \|v\|\), para todo \(\lambda \in \mathbb{F}\).
    \end{enumerate}
\end{prop}

\begin{defn}
    Dois vetores \(u, v \in V\) são ditos \emph{ortogonais} se \(\langle u, v \rangle = 0\).
\end{defn}

\begin{prop} (ortogonalidade e o zero).
    \begin{enumerate}
        \item [\textup{(a)}] O vetor nulo é ortogonal a todo vetor em \(V\).
        \item [\textup{(b)}] O vetor nulo é o único vetor em \(V\) que é ortogonal a si mesmo.
    \end{enumerate}
\end{prop}

\begin{thm} (teorema de Pitágoras).
    Sejam \(u, v \in V\). Se \(u\) e \(v\) são ortogonais, então
    \[
        \|u + v\|^2 = \|u\|^2 + \|v\|^2.
    \]
\end{thm}

\begin{prop} (decomposição ortogonal).
    Suponha que \(u, v \in V\), com \(v\neq0\). Faça \(\displaystyle c=\frac{\langle u,v \rangle}{\Vert v \Vert^2}\) e \(\displaystyle w = u - \frac{\langle u,v \rangle}{\Vert v \Vert^2} v\). Então \(u = cv + w\) e \(\langle w, v \rangle = 0\).
\end{prop}

\begin{thm} (desigualdade de Cauchy-Schwarz).
    Sejam \(u, v \in V\). Então
    \begin{equation}
        |\langle u, v \rangle| \leq \|u\| \|v\|.
    \end{equation}
    \noindent Além disso, a igualdade ocorre sse um de \(u, v\) é um múltiplo escalar do outro.
\end{thm}

\begin{cor} (desigualdade triangular).
    Sejam \(u, v \in V\). Então
    \begin{equation}
        \|u + v\| \leq \|u\| + \|v\|.
    \end{equation}
    \noindent Essa desigualdade é uma igualdade sse um de \(u, v\) é um múltiplo escalar do outro com escalar não-negativo.
\end{cor}

\begin{cor} (igualdade do paralelogramo).
    Sejam \(u, v \in V\). Então
    \begin{equation}
        \|u + v\|^2 + \|u - v\|^2 = 2\|u\|^2 + 2\|v\|^2.
    \end{equation}
\end{cor}

\subsection{Bases Ortonormais}

\begin{defn} (lista ortonormal).
    \begin{itemize}
        \item Uma lista de vetores é dita \emph{ortonormal} se cada vetor na lista tem norma \(1\) e é ortogonal a todos os outros vetores na lista.
        \item Em outras palavras, uma lista \(e_1, \ldots, e_m\) de vetores em \(V\) é ortonormal se
        \[
            \langle e_j, e_k \rangle = \begin{cases}
                1 & \text{se } j = k \\
                0 & \text{se } j \neq k
            \end{cases}, \quad \forall j,k = 1,\ldots,m. 
        \]
    \end{itemize}
\end{defn}

\begin{prop} (norma de uma combinação linear de vetores ortonormais).
    Se \(e_1, \ldots, e_m\) é uma lista ortonormal de vetores em \(V\), então
    \begin{equation}
        \|a_1 e_1 + a_2 e_2 + \cdots + a_m e_m\|^2 = |a_1|^2 + |a_2|^2 + \cdots + |a_m|^2,
    \end{equation}
    \noindent para todos os escalares \(a_1, a_2, \ldots, a_m \in \mathbb{F}\).
\end{prop}

\begin{cor} (listas ortonormais são linearmente independentes).
    Toda lista ortonormal de vetores é linearmente independente.
\end{cor}

\begin{prop} (desigualdade de Bessel).
    Suponha que \(e_1, \ldots, e_m\) é uma lista ortonormal de vetores em \(V\). Então, para todo \(v \in V\),
    \begin{equation}
        |\langle v, e_1 \rangle|^2 + |\langle v, e_2 \rangle|^2 + \cdots + |\langle v, e_m \rangle|^2 \leq \|v\|^2.
    \end{equation}
\end{prop}

\begin{defn}
    Uma \emph{base ortonormal} de \(V\) é uma base de \(V\) que é uma lista ortonormal de vetores.
\end{defn}

\begin{prop}
    Seja \(V\) um espaço vetorial de dimensão finita com produto interno. Então toda lista ortonormal de vetores em \(V\) de tamanho \(\dim V\) é uma base ortonormal de \(V\).
\end{prop}

\begin{prop} (vetor como combinação linear de uma base ortonormal).
    Suponha que \(e_1, \ldots, e_n\) é uma base ortonormal de \(V\). Então, para todo \(v \in V\),
    \begin{enumerate}
        \item [\textup{(a)}] \(v = \langle v, e_1 \rangle e_1 + \langle v, e_2 \rangle e_2 + \cdots + \langle v, e_n \rangle e_n\) (identidade de Parseval);
        \item [\textup{(b)}] \(\|v\|^2 = |\langle v, e_1 \rangle|^2 + |\langle v, e_2 \rangle|^2 + \cdots + |\langle v, e_n \rangle|^2\).
        \item [\textup{(c)}] \(\langle u, v \rangle = \langle u, e_1 \rangle \overline{\langle v, e_1 \rangle} + \langle u, e_2 \rangle \overline{\langle v, e_2 \rangle} + \cdots + \langle u, e_n \rangle \overline{\langle v, e_n \rangle}\), para todo \(u \in V\).
    \end{enumerate}
\end{prop}

\begin{thm} (processo de ortogonalização de Gram-Schmidt).
    Suponha que \(v_1, \ldots, v_m\) é uma lista de vetores linearmente independente em \(V\). Faça \(f_1 = v_1\). Para \(k=2, \ldots, m\), defina \(f_k\) indutivamente por 
    \[
        f_k = v_k - \sum_{j=1}^{k-1} \frac{\langle v_k, f_j \rangle}{\|f_j\|^2} f_j.
    \]
    \noindent Para todo \(k = 1, \ldots, m\), defina \(e_k = \frac{f_k}{\|f_k\|}\). Então \(e_1, \ldots, e_m\) é uma lista ortonormal de vetores em \(V\) tal que
    \[
        \spn(e_1, \ldots, e_k) = \spn(v_1, \ldots, v_k),
    \]
    \noindent para todo \(k = 1, \ldots, m\).
\end{thm}

\begin{prop} (existência de base ortonormal).
    Todo espaço vetorial de dimensão finita com produto interno tem uma base ortonormal.
\end{prop}

\begin{cor} (toda lista ortonormal pode ser estendida a uma base ortonormal).
    Suponha que \(V\) é um espaço vetorial de dimensão finita com produto interno. Então toda lista ortonormal de vetores em \(V\) pode ser estendida a uma base ortonormal de \(V\).
\end{cor}

\begin{prop} (matriz triangular-superior com respeito a uma base ortonormal).
    Seja \(V\) um espaço vetorial de dimensão finita com produto interno e \(T \in \mathcal{L}(V)\). Então \(T\) tem uma matriz triangular-superior com respeito a alguma base ortonormal de \(V\) sse o polinômio mínimo de \(T\) é igual a \((z - \lambda_1) \cdots (z - \lambda_m)\), para alguns \(\lambda_1, \ldots, \lambda_m \in \mathbb{F}\).
\end{prop}

\begin{thm} (teorema de Schur).
    Todo operador em um espaço vetorial complexo de dimensão finita tem uma matriz triangular-superior com respeito a alguma base ortonormal.
\end{thm}

\begin{thm} (teorema de representação de Riesz).
    Suponha que \(V\) é um espaço vetorial de dimensão finita com produto interno e \(\varphi \in V'\). Então existe um único vetor \(v \in V\) tal que
    \[
        \varphi(u) = \langle u, v \rangle,
    \]
    para todo \(u \in V\).
\end{thm}

\subsection{Complementos Ortogonais e Problemas de Minimização}

\begin{defn} (complemento ortogonal, \(U^\perp\)).
    Se \(U\) é um subconjunto de \(V\), então o \emph{complemento ortogonal} de \(U\), denotado por \(U^\perp\), é o conjunto de todos os vetores em \(V\) que são ortogonais a todos os vetores em \(U\):
    \[
        U^\perp = \{v \in V : \langle v, u \rangle = 0 \quad \forall\,  u \in U\}.
    \]
\end{defn}

\begin{prop} (propriedades do complemento ortogonal).
    \begin{enumerate}
        \item [\textup{(a)}] Se \(U\) é um subconjunto de \(V\), então \(U^\perp\) é um subespaço de \(V\).
        \item [\textup{(b)}] \(\{0\}^\perp = V\).
        \item [\textup{(c)}] \(V^\perp = \{0\}\).
        \item [\textup{(d)}] Se \(U\) é um subconjunto de \(V\), então \(U \cap U^\perp \subset \{0\}\).
        \item [\textup{(e)}] Se \(U\) e \(W\) são subconjuntos de \(V\) com \(U \subset W\), então \(W^\perp \subset U^\perp\).
    \end{enumerate}
\end{prop}

\begin{prop} (soma direta de um subespaço e seu complemento ortogonal).
    Seja \(U\) um subespaço de um espaço vetorial de dimensão finita com produto interno \(V\). Então
    \[
        V = U \oplus U^\perp.
    \]
\end{prop}

\begin{cor} (dimensão do complemento ortogonal).
    Seja \(U\) um subespaço de um espaço vetorial de dimensão finita com produto interno \(V\). Então
    \[
        \dim U^\perp = \dim V - \dim U.
    \]
\end{cor}

\begin{cor} (complemento ortogonal do complemento ortogonal).
    Seja \(U\) um subespaço de um espaço vetorial de dimensão finita com produto interno \(V\). Então
    \[
        (U^\perp)^\perp = U.
    \]
\end{cor}

\begin{cor} (\(U^\perp = \{0\} \iff U = V\) para \(U\) subespaço de dimensão finita de \(V\)).
    Suponha que \(U\) é um subespaço de um espaço vetorial de dimensão finita com produto interno \(V\). Então 
    \[
        U^\perp = \{0\} \iff U = V.
    \]
\end{cor}

\begin{defn} (projeção ortogonal).
    Seja \(U\) um subespaço de um espaço vetorial de dimensão finita com produto interno \(V\). A \emph{projeção ortogonal} de \(V\) sobre \(U\) é o operador \(P_U \in \mathcal{L}(V)\) definido por 
    \[
        P_U v = u,
    \]
    \noindent onde \(v = u + w\) com \(u \in U\) e \(w \in U^\perp\), para todo \(v \in V\).
\end{defn}

\begin{prop} (propriedades da projeção ortogonal).
    Suponha que \(U\) é um subespaço de um espaço vetorial de dimensão finita com produto interno \(V\) e \(P_U\) é a projeção ortogonal de \(V\) sobre \(U\). Então
    \begin{enumerate}
        \item [\textup{(a)}] \(P_U \in \mathcal{L}(V)\);
        \item [\textup{(b)}] \(P_U u = u\), para todo \(u \in U\);
        \item [\textup{(c)}] \(P_U w = 0\), para todo \(w \in U^\perp\);
        \item [\textup{(d)}] \(\im P_U = U\);
        \item [\textup{(e)}] \(\ker P_U = U^\perp\);
        \item [\textup{(f)}] \(v - P_U v \in U^\perp\), para todo \(v \in V\).
        \item [\textup{(g)}] \(P_U^2 = P_U\);
        \item [\textup{(h)}] \(\|P_U v\| \leq \|v\|\), para todo \(v \in V\).
        \item [\textup{(i)}] se \(e_1, \ldots, e_m\) é uma base ortonormal de \(U\) e \(v \in V\), então
        \[
            P_U v = \langle v, e_1 \rangle e_1 + \langle v, e_2 \rangle e_2 + \cdots + \langle v, e_m \rangle e_m.
        \]
    \end{enumerate}
\end{prop}

\begin{prop} (minimização da distância a um subespaço).
    Seja \(U\) um subespaço de um espaço vetorial de dimensão finita com produto interno \(V\), \(v \in V\) e \(u \in U\). Então 
    \begin{equation}
        \| v - P_U v \| \leq \| v - u \|.
    \end{equation}
    Além disso, a igualdade ocorre sse \(u = P_U v\).
\end{prop}

\section{Operadores em Espaços com Produto Interno}

\subsection{Operadores Autoadjuntos e Operadores Normais}

\begin{defn} (adjunta, \(T^*\)).
    Suponha que \(T \in \mathcal{L}(V, W)\). A \emph{adjunta} de \(T\) é a função \(T^*: W \to V\) tal que

    \begin{equation} \label{eq:adjunta}
        \langle Tv, w \rangle = \langle v, T^* w \rangle,
    \end{equation}

    \noindent para todo \(v \in V\) e todo \(w \in W\).
\end{defn}

\begin{rem}
    Na equação \eqref{eq:adjunta}, o produto interno à esquerda é o produto interno em \(W\) e o produto interno à direita é o produto interno em \(V\).
\end{rem}

\begin{prop} (adjunta de uma transformação linear é uma transformação linear).
    Se \(T \in \mathcal{L}(V, W)\), então \(T^* \in \mathcal{L}(W, V)\).
\end{prop}

\begin{prop} (propriedades da adjunta).
    \begin{enumerate}
        \item [\textup{(a)}] \((S + T)^* = S^* + T^*\), para todos \(S, T \in \mathcal{L}(V, W)\).
        \item [\textup{(b)}] \((\lambda T)^* = \overline{\lambda} T^*\), para todo \(\lambda \in \mathbb{F}\) e todo \(T \in \mathcal{L}(V, W)\).
        \item [\textup{(c)}] \((T^*)^* = T\);
        \item [\textup{(d)}] \((ST)^* = T^* S^*\), para todos \(S \in \mathcal{L}(W, X)\) e \(T \in \mathcal{L}(V, W)\).
        \item [\textup{(e)}] \(I^* = I\), onde \(I\) é o operador identidade em \(V\).
        \item [\textup{(f)}] Se \(T\) é invertível, então \(T^*\) é invertível e \((T^*)^{-1} = (T^{-1})^*\).
    \end{enumerate}
\end{prop}

\begin{prop} (espaço nulo e imagem da adjunta).
    Se \(T \in \mathcal{L}(V, W)\), então
    \begin{enumerate}
        \item [\textup{(a)}] \(\ker T^* = (\im T)^\perp\);
        \item [\textup{(b)}] \(\im T^* = (\ker T)^\perp\);
        \item [\textup{(c)}] \(\ker T = (\im T^*)^\perp\);
        \item [\textup{(d)}] \(\im T = (\ker T^*)^\perp\).
    \end{enumerate}
\end{prop}

\begin{defn} (conjugada transposta, \(A^*\)).
    A \emph{conjugada tranposta} de uma matriz \(A\) \(m\) por \(n\) é a matriz \(A^*\) \(n\) por \(m\) definida por
    \[
        (A^*)_{jk} = \overline{A_{kj}},
    \]
    \noindent para todo \(j = 1, \ldots, n\) e todo \(k = 1, \ldots, m\).
\end{defn}

\begin{prop} (matriz da adjunta é igual à conjugada transposta da matriz).
    Seja \(T \in \mathcal{L}(V, W)\). Se \(e_1, \ldots, e_n\) é uma base ortonormal de \(V\) e \(f_1, \ldots, f_m\) é uma base ortonormal de \(W\), então
    \[
        \mathcal{M}(T^*, (f_1, \ldots, f_m), (e_1, \ldots, e_n)) = (\mathcal{M}(T, (e_1, \ldots, e_n), (f_1, \ldots, f_m)))^*.
    \]
\end{prop}

\begin{defn} (operador autoadjunto).
    Um operador \(T \in \mathcal{L}(V)\) é chamado de \emph{autoadjunto} se \(T = T^*\).
\end{defn}

\begin{rem}
    Um operador \(T \in \mathcal{L}(V)\) é autoadjunto sse
    \[
        \langle Tv, w \rangle = \langle v, Tw \rangle
    \]
    \noindent para todos \(v, w \in V\).
\end{rem}

\begin{prop} (autovalores de operadores autoadjuntos).
    Todo autovalor de um operador autoadjunto é real.
\end{prop}

\begin{prop} (sobre \(\mathbb{C}\), \(Tv\) é ortogonal a \(v\) para todo \(v \iff T = 0\)).
    Suponha que \(V\) é um espaço vetorial complexo com produto interno e \(T \in \mathcal{L}(V)\). Então \(Tv\) é ortogonal a \(v\) para todo \(v \in V\) sse \(T = 0\).
\end{prop}

\begin{prop} (sobre \(\mathbb{C}\), \(\langle Tv, v \rangle\) é real para todo \(v \iff T\) é autoadjunto).
    Suponha que \(V\) é um espaço vetorial complexo com produto interno e \(T \in \mathcal{L}(V)\). Então \(\langle Tv, v \rangle\) é real para todo \(v \in V\) sse \(T\) é autoadjunto.
\end{prop}

\begin{prop} (\(T\) é autoadjunto e \(\langle Tv, v \rangle = 0\) para todo \(v \iff T = 0\)).
    Suponha que \(T\) é um operador autoadjunto em \(V\). Então \(\langle Tv, v \rangle = 0\) para todo \(v \in V\) sse \(T = 0\).
\end{prop}

\begin{defn} (operador normal).
    Um operador \(T \in \mathcal{L}(V)\) é chamado de \emph{normal} se \(T T^* = T^* T\).
\end{defn}

\begin{rem}
    Todo operador autoadjunto é normal.
\end{rem}

\begin{prop} (\(T\) é normal sse \(Tv\) e \(T^* v\) têm a mesma norma).
    Suponha que \(T \in \mathcal{L}(V)\). Então \(T\) é normal se, e somente se, \(\|Tv\| = \|T^* v\|\) para todo \(v \in V\).
\end{prop}

\begin{prop} (imagem, núcleo e autovetores de um operador normal).
    Seja \(T \in \mathcal{L}(V)\) um operador normal. Então
    \begin{enumerate}
        \item [\textup{(a)}] \(\ker T = \ker T^*\);
        \item [\textup{(b)}] \(\im T = \im T^*\);
        \item [\textup{(c)}] \(V = \ker T \oplus \im T\);
        \item [\textup{(d)}] \(T - \lambda I\) é normal, para todo \(\lambda \in \mathbb{F}\);
        \item [\textup{(e)}] se \(v \in V\) e \(\lambda \in \mathbb{F}\), então \(Tv = \lambda v \iff T^* v = \overline{\lambda} v\).
    \end{enumerate}
\end{prop}

\begin{prop} (autovetores ortogonais para operadores normais).
    Suponha que \(T \in \mathcal{L}(V)\) é um operador normal. Então autovetores de \(T\) correspondentes a autovalores distintos são ortogonais.
\end{prop}

\begin{prop} (\(T\) é normal sse as partes real e imaginária de \(T\) comutam).
    Seja \(\mathbb{F} = \mathbb{C}\) e \(T \in \mathcal{L}(V)\). Então \(T\) é normal sse existem operadores autoadjuntos \(A, B \in \mathcal{L}(V)\) tais que \(T = A + iB\) e \(AB = BA\).
\end{prop}

\subsection{Teorema Espectral}

\begin{thm} (teorema espectral sobre \(\mathbb{R}\)).
    Suponha que \(V\) é um espaço vetorial real de dimensão finita com produto interno e \(T \in \mathcal{L}(V)\). Então as seguintes condições são equivalentes.
    \begin{enumerate}
        \item[\textup{(a)}] \(T\) é autoadjunto.
        \item[\textup{(b)}] \(V\) tem uma base ortonormal de autovetores de \(T\).
        \item[\textup{(c)}] \(T\) tem uma matriz diagonal com respeito a alguma base ortonormal de \(V\).
    \end{enumerate}
\end{thm}

\begin{thm} (teorema espectral sobre \(\mathbb{C}\)).
    Suponha que \(V\) é um espaço vetorial complexo de dimensão finita com produto interno e \(T \in \mathcal{L}(V)\). Então as seguintes condições são equivalentes.
    \begin{enumerate}
        \item[\textup{(a)}] \(T\) é normal.
        \item[\textup{(b)}] \(V\) tem uma base ortonormal de autovetores de \(T\).
        \item[\textup{(c)}] \(T\) tem uma matriz diagonal com respeito a alguma base ortonormal de \(V\).
    \end{enumerate}
    
\end{thm}

\subsection{Operadores Positivos}

\begin{defn} (operador positivo).
    Um operador \(T \in \mathcal{L}(V)\) é chamado de \emph{positivo} se \(T\) é autoadjunto e \(\langle Tv, v \rangle \geq 0\) para todo \(v \in V\).
\end{defn}

\begin{defn}
    Um operador \(R\) é chamado de \emph{raiz quadrada} de um operador \(T\) se \(R^2 = T\).
\end{defn}

\subsection{Isometrias, Operadores Unitários e Fatoração Matricial}
\subsection{Decomposição em Valores Singulares}
\subsection{Consequências de Valores Singulares}

\section{Operadores em Espaços Vetoriais Complexos}
\section{Álgebra Multilinear e Determinantes}

\end{document}